% Saloum poster (French). Build: cd poster && lualatex poster_saloum_fr.tex
% Figures are produced by the scripts in ../scripts (see README). UCAD logo: Wikimedia Commons, CC BY-SA 4.0.
% Poster scientifique -- mangroves du Saloum, NISAR (bande L) et Sentinel-1 (bande C)
\documentclass[20pt,a0paper,portrait,margin=0mm,innermargin=18mm,blockverticalspace=14mm,colspace=16mm,subcolspace=8mm]{tikzposter}
\usepackage[french]{babel}
\usepackage{fontspec}
\setmainfont{TeX Gyre Heros}
\usepackage{graphicx,booktabs,xcolor,amsmath,enumitem,tabularx}
\usetikzlibrary{positioning,arrows.meta,fit,calc}
\usepackage[hidelinks]{hyperref}
\definecolor{ucad}{HTML}{2A91CF}
\definecolor{nuit}{HTML}{12304F}
\definecolor{bandeL}{HTML}{1F5FAE}
\definecolor{bandeC}{HTML}{D4541F}
\definecolor{fond}{HTML}{F5F8FB}
\definecolor{gris}{HTML}{555555}
\usetheme{Simple}
\colorlet{backgroundcolor}{white}
\colorlet{framecolor}{white}
\colorlet{blocktitlebgcolor}{nuit}
\colorlet{blocktitlefgcolor}{white}
\colorlet{blockbodybgcolor}{fond}
\colorlet{blockbodyfgcolor}{black}
\colorlet{titlebgcolor}{white}
\colorlet{titlefgcolor}{nuit}
\defineblockstyle{Plat}{titlewidthscale=1, bodywidthscale=1, titleleft, titleoffsetx=0pt, titleoffsety=0pt,
  bodyoffsetx=0mm, bodyoffsety=12mm, bodyverticalshift=12mm, roundedcorners=6, linewidth=0pt,
  titleinnersep=9mm, bodyinnersep=11mm}{
  \draw[color=framecolor, fill=blockbodybgcolor, rounded corners=\blockroundedcorners] (blockbody.south west) rectangle (blockbody.north east);
  \ifBlockHasTitle
    \draw[color=framecolor, fill=blocktitlebgcolor, rounded corners=\blockroundedcorners] (blocktitle.south west) rectangle (blocktitle.north east);
  \fi}
\useblockstyle{Plat}
\setlist{leftmargin=1.1em,itemsep=4pt,topsep=4pt}
\newcommand{\Lb}[1]{\textcolor{bandeL}{\textbf{#1}}}
\newcommand{\Cb}[1]{\textcolor{bandeC}{\textbf{#1}}}
\newcommand{\fig}[2][1]{\begin{center}\includegraphics[width=#1\linewidth]{#2}\end{center}}
\newcommand{\legende}[1]{{\small\color{gris}#1\par}}
\newcounter{numfig}\newcounter{numtab}
\newcommand{\figcap}[1]{\stepcounter{numfig}{\small\textbf{Figure \thenumfig.} #1\par}\vspace{2mm}}
\newcommand{\tabcap}[1]{\stepcounter{numtab}{\small\textbf{Tableau \thenumtab.} #1\par}\vspace{1mm}}
\newcommand{\sous}[1]{\par\vspace{3mm}{\color{nuit}\large\bfseries #1\par}\vspace{1mm}}
\newcommand{\cle}[1]{\vspace{4mm}\colorbox{ucad!15}{\parbox{0.97\linewidth}{\vspace{3mm}#1\vspace{3mm}}}\vspace{2mm}}

\settitle{
\newcommand{\titrebloc}{\centering
{\color{nuit}\fontsize{66}{78}\selectfont\bfseries Les mangroves du Saloum vues par NISAR\\[5mm]
\fontsize{50}{60}\selectfont bande L contre bande C, et la marée au passage du satellite\par}
\vspace{9mm}
{\fontsize{40}{48}\selectfont\color{black}\textbf{Mamadou Galy DIALLO}\par}
\vspace{3mm}
{\fontsize{30}{38}\selectfont\color{gris}galy.quantum@gmail.com\par}}
\ifdefined\sansLogo
\begin{minipage}[c]{\linewidth}\titrebloc\end{minipage}
\else
\begin{minipage}[c]{0.14\linewidth}\centering\includegraphics[width=0.72\linewidth]{logos/logo_ucad.png}\end{minipage}%
\begin{minipage}[c]{0.72\linewidth}\titrebloc\end{minipage}%
\begin{minipage}[c]{0.14\linewidth}\centering\includegraphics[width=0.72\linewidth]{logos/logo_ucad.png}\end{minipage}
\fi
\vspace{8mm}
{\color{ucad}\rule{\linewidth}{5mm}}
}

\tikzposterlatexaffectionproofoff
\graphicspath{{../}}   % figures/ at the repository root
\begin{document}
\maketitle[width=0.95\paperwidth, titletotopverticalspace=12mm, titletoblockverticalspace=10mm]

\begin{columns}
% ===================================================================== colonne gauche
\column{0.5}

\block{1. Introduction}{
Les mangroves comptent parmi les écosystèmes côtiers les plus riches en carbone [3] et protègent les
littoraux ; les suivre régulièrement est donc essentiel dans des milieux qui changent vite. Au Sénégal, les
mangroves du delta du Saloum, longtemps décrites comme en recul, s'étendent de nouveau depuis les
années 2000 [4] : il faut pouvoir mesurer ces gains comme les pertes. Le radar en \Lb{bande L}, qui traverse
le houppier et ne dépend ni des nuages ni de l'éclairement, est la référence pour cartographier les
mangroves [1, 2]. La détection des changements reste pourtant son point faible : 58 à 61\,\% d'exactitude
seulement pour les gains et pertes de Global Mangrove Watch [2]. Une source d'erreur possible est l'eau :
sous le couvert, la hauteur de la marée au moment du passage modifie le signal radar ; dans les marais
littoraux, la rétrodiffusion L baisse quand l'eau monte [5]. La marée est donc le plus souvent traitée
comme un bruit, et les travaux sur le Saloum reposent encore sur l'optique ou sur la \Cb{bande C} [4].
Depuis 2025, la mission \textbf{NISAR} (NASA--ISRO) fournit pour la première fois une série L gratuite tous
les \textbf{12 jours}, ce qui permet de poser ces questions depuis l'espace. \textbf{Ainsi, cette étude
propose une approche comparative de la bande L (NISAR) et de la bande C (Sentinel-1) sur les mangroves du
delta du Saloum.}
\sous{Objectif}
L'étude s'appuie sur onze dates NISAR et seize dates Sentinel-1 acquises de juin à novembre 2026, en intensité
et en cohérence, et sur la
marée observée au marégraphe de Dakar. Elle cherche d'abord à savoir laquelle des deux bandes distingue le
mieux la mangrove des autres milieux du delta. Elle vérifie ensuite que les passages de NISAR, toujours à la
même heure, échantillonnent bien tout le cycle de marée, puis estime jusqu'où la marée peut modifier la
rétrodiffusion et la cohérence de la mangrove en bande L.
}

\block{2. Zone d'étude}{
\fig{figures/f0_localisation.png}
\figcap{Zone d'étude. a) Localisation : Sénégal (limites GADM 4.1), fenêtre d'étude et emprises réelles des
acquisitions radar (NISAR trace 124 D, Sentinel-1 orbite 133 A). b) Composition colorée NISAR, bande L
(R = HH, V = HV, B = HH/HV). c) Composition colorée Sentinel-1, bande C (R = VV, V = VH, B = VV/VH).
Médianes temporelles juin--novembre 2026 ; l'eau apparaît en noir.}
Le \textbf{delta du Saloum} (fenêtre de 33 $\times$ 33 km, réserve de biosphère UNESCO) est un estuaire
inverse, hypersalin, où des mangroves basses à \textit{Avicennia} et \textit{Rhizophora} bordent des chenaux
de marée (bolongs) et des étendues nues salées (tannes). La mangrove couvre environ 414 km$^2$ de la
fenêtre (WorldCover). Marégraphe de référence : Dakar, à environ 120 km du centre de la fenêtre.
}

\block{3. Données et méthodes}{
\tabcap{Données (juin -- novembre 2026, toutes en accès libre).}
\begin{center}\small
\begin{tabularx}{\linewidth}{@{}lXcc@{}}
\toprule
Produit & Variables & Géométrie & Acquisitions \\
\midrule
\Lb{NISAR L2 GCOV} & $\gamma^0$ HH, HV & trace 124 D, 19 h 34 UTC & 11 dates \\
\Lb{NISAR L2 GUNW} & cohérence HH à 12 jours & trace 124 D & 8 paires \\
\Cb{Sentinel-1 RTC} & $\gamma^0$ VV, VH & orbite 133 A, 19 h 17 UTC & 16 dates \\
\Cb{Sentinel-1 SLC} & cohérence VV à 12 jours (SNAP) & orbite 133 A & 13 paires \\
ESA WorldCover 2021 & classes d'occupation du sol & -- & 10 m \\
Marégraphe de Dakar & hauteur d'eau & -- & UHSLC 2015--22, IOC 2026 \\
\bottomrule
\end{tabularx}\end{center}
\vspace{4mm}
\begin{center}\resizebox{\linewidth}{!}{%
\begin{tikzpicture}[font=\normalsize, node distance=9mm and 7mm,
  b/.style={draw=#1, fill=#1!10, rounded corners=3mm, line width=1.2mm, align=center, minimum height=26mm, text width=#1w},
  bx/.style={draw=nuit!70, fill=white, rounded corners=3mm, line width=1mm, align=center, minimum height=24mm},
  fl/.style={-{Latex[length=7mm,width=6mm]}, line width=1.4mm, nuit!70}]
\node[bx, draw=bandeL, fill=bandeL!8, text width=11cm] (gcov) {\textbf{NISAR GCOV + GUNW}\\ $\gamma^0$ HH/HV, cohérence};
\node[bx, draw=bandeC, fill=bandeC!8, text width=9cm, right=of gcov] (s1) {\textbf{Sentinel-1 RTC + SLC}\\ $\gamma^0$ VV/VH, cohérence};
\node[bx, text width=8cm, right=of s1] (wc) {\textbf{WorldCover}\\ classes};
\node[bx, draw=green!50!black, fill=green!6, text width=12cm, right=of wc] (tide) {\textbf{Marégraphes}\\ UHSLC 2015--22, IOC 2026};
\node[bx, text width=30cm, below=14mm of $(gcov.south)!0.5!(wc.south)$] (grid) {\textbf{Paquet \texttt{nisar\_s1}} : recherche ASF, lecture à distance par fenêtre,\\ grille commune UTM 20 m, médianes temporelles par pixel, classes érodées d'un pixel};
\node[bx, text width=12cm, below=14mm of tide] (model) {\textbf{Modèle harmonique}\\ 22 constituantes ; hauteur à chaque passage, \textbf{contrôlée sur IOC}};
\node[bx, fill=ucad!12, text width=19cm, below=14mm of grid.south west, anchor=north west] (sep) {\textbf{Séparabilité de la mangrove}\\ AUC mangrove / classe, IC 95\,\% par bootstrap par blocs (500 m)};
\node[bx, fill=ucad!12, text width=21cm, right=of sep] (sens) {\textbf{Sensibilité à la marée}\\ $\gamma^0$ = f(marée) : pente et IC 95\,\% ; cohérence = f($|\Delta h|$)};
\draw[fl] (gcov.south) -- (gcov.south|-grid.north);
\draw[fl] (s1.south) -- (s1.south|-grid.north);
\draw[fl] (wc.south) -- (wc.south|-grid.north);
\draw[fl] (tide) -- (model);
\draw[fl] (grid.south-|sep.north) -- (sep.north);
\draw[fl] ($(sens.north west)+(5cm,0)$|-grid.south) -- ($(sens.north west)+(5cm,0)$);
\draw[fl] (model.south) -- (model.south|-sens.north);
\end{tikzpicture}}
\end{center}
\figcap{Schéma de la méthode. Toute la chaîne radar est dans le paquet libre \texttt{nisar\_s1}. AUC : aire sous la courbe ROC, soit la probabilité qu'un pixel de mangrove et un pixel de l'autre classe tirés au hasard soient correctement ordonnés (0,5 = indiscernables, 1 = séparation parfaite).}
}

% ===================================================================== colonne droite
\column{0.5}

\block{4. Résultats}{
\sous{4.1 La bande L isole la mangrove de la forêt de terre ferme}
\fig{figures/f2_separabilite.png}
\figcap{AUC de séparation (axe horizontal) entre la mangrove et chaque classe d'occupation du sol WorldCover, pour l'intensité (L HV, L HH, C VH, C VV) et la cohérence à 12 jours (L et C) ; 0,5 = aucune séparation, 1 = séparation parfaite ; tirets : AUC = 0,8.}
\cle{Contre la \textbf{forêt continentale} : AUC \Lb{0,87 en L HV} [0,86--0,89] contre \Cb{0,54 en C VH}
[0,51--0,57] ; en cohérence, \Lb{0,77 en L} [0,75--0,79] contre \Cb{0,63 en C} [0,62--0,65] (IC 95\,\%).
Contre l'eau, les tannes, les cultures et la savane, les intensités des deux bandes séparent bien
(AUC $\geq$ 0,9) ; contre les arbustes, aucune variable ne suffit (AUC $\leq$ 0,71).}

\sous{4.2 NISAR échantillonne tout le cycle de marée, vérifié sur mesures}
\fig{figures/f3_maree.png}
\figcap{Marée prédite à Dakar par le modèle harmonique (courbe grise), hauteur prédite à chaque passage
Sentinel-1 (ronds) et NISAR (losanges pleins), et hauteur mesurée au marégraphe IOC aux passages NISAR
(losanges vides), juin--novembre 2026.}
La phase M2 avance d'environ 68$^\circ$ en 12 jours : les 11 passages couvrent \textbf{1,00 m} de marée
prédite. Le modèle s'écarte des mesures IOC de \textbf{4,5 cm} aux passages NISAR (7, $\rho$ = 0,89) et de
5,0 cm aux passages Sentinel-1 (11) ; les passages mesurés couvrent \textbf{0,90 m}.

\sous{4.3 La mangrove varie peu avec la marée, en bande L comme en bande C}
\fig{figures/f4_stabilite.png}
\figcap{Médiane de $\gamma^0$ de la mangrove, en écart à sa moyenne, en bandes L (HV, HH) et C (VH, VV) : a) à chaque date de juin à novembre 2026 ; b) en fonction de la hauteur de marée au passage du satellite.}
\vspace{2mm}
\tabcap{Pente de $\gamma^0$ de la mangrove en fonction de la marée au passage (régression linéaire, valeurs brutes). Aucune pente n'est significative : tous les IC contiennent 0.}
\begin{center}{\renewcommand{\arraystretch}{1.0}\setlength{\tabcolsep}{14pt}
\begin{tabular}{lccccc}
\toprule
 & dates & amplitude & pente & IC 95\,\% & $p$ \\
 &  & (dB) & (dB/m) & (dB/m) & \\
\midrule
\Lb{L HV} & 11 & 0,52 & +0,13 & $-$0,23 à +0,48 & 0,44 \\
\Lb{L HH} & 11 & 0,95 & +0,36 & $-$0,18 à +0,91 & 0,17 \\
\Cb{C VH} & 16 & 0,77 & +0,22 & $-$0,20 à +0,63 & 0,28 \\
\Cb{C VV} & 16 & 0,93 & $-$0,15 & $-$0,70 à +0,40 & 0,56 \\
\bottomrule
\end{tabular}}\end{center}


\sous{4.4 Cohérence : élevée en bande L, effondrée en bande C, indépendante de la marée}
\fig{figures/f5_coherence.png}
\figcap{Cohérence à 12 jours. a, b) Médiane temporelle sur la fenêtre d'étude : a) NISAR, bande L, HH
(8 paires) ; b) Sentinel-1, bande C, VV (13 paires) ; eau permanente en gris. c, d) Médiane par classe en
fonction de l'écart de marée $|\Delta h|$ entre les deux dates : c) bande L (mangrove $\rho$ = +0,26,
$p$ = 0,53) ; d) bande C (mangrove $\rho$ = +0,05, $p$ = 0,87).}
En bande L, les îles de mangrove ressortent (médiane \Lb{0,72}) ; en bande C, elles sont presque
uniformément décorrélées (\Cb{0,27}) : le houppier décorrèle la bande C, et seuls les tannes et le bâti
y restent cohérents. Aucune classe ne varie significativement avec $|\Delta h|$ (toutes les $p \geq 0{,}07$).
}

\end{columns}

\begin{columns}
\column{0.67}
\block{5. Discussion et conclusions}{
Cette étude propose une approche comparative des bandes L (NISAR) et C (Sentinel-1) pour distinguer les
mangroves du Saloum et évaluer l'effet de la marée au passage. En intensité comme en cohérence, la bande
L se démarque face à la \textbf{forêt de terre ferme} : l'AUC passe de \Cb{0,54} (C VH) à \Lb{0,87}
(L HV), et de \Cb{0,63} à \Lb{0,77} en cohérence. La bande L est sensible à la structure et à la
biomasse [1] ; en bande C, le signal vient des feuilles du houppier, qui bougent, et la cohérence de la
mangrove tombe à \Cb{0,27} contre \Lb{0,72}. Les passages NISAR couvrent \textbf{0,90 m} de marée
mesurée, presque tout le marnage, mais la mangrove varie de moins de \textbf{0,5 dB} en L HV, sans pente
significative, ce qui rejoint les travaux sur les marais littoraux [5]. L'effet reste borné plutôt
qu'exclu : une seule saison des pluies, une marée mesurée à Dakar et non dans les bolongs, une référence
WorldCover non vérifiée sur le terrain. \textbf{Ainsi, en perspective}, la saison sèche, la marée du
delta et la validation de terrain sont des pistes à explorer, en préparation de la mission \textbf{ROSE-L}.
}
\column{0.33}
\block{Références}{\footnotesize
\begin{enumerate}[label={[\arabic*]},leftmargin=2.1em,itemsep=2pt]
\item Lucas \textit{et al.} (2007), \textit{Aquatic Conservation} 17, 245--264.
\item Bunting \textit{et al.} (2022), \textit{Remote Sensing} 14, 3657.
\item Simard \textit{et al.} (2019), \textit{Nature Geoscience} 12, 40--45.
\item Lombard \& Andrieu (2021), \textit{Remote Sensing} 13, 1961.
\item Lamb \textit{et al.} (2025), \textit{Remote Sensing} 17, 263.
\end{enumerate}
\vspace{3mm}
{\color{gris}\textbf{Données} : NISAR L2 (NASA--ISRO ; produits JPL générés avec \textbf{ISCE3}, ASF DAAC),
Copernicus Sentinel-1 (ESA, UE), ESA WorldCover 2021, marégraphes UHSLC et IOC (VLIZ/UNESCO).
\textbf{Code} : \texttt{nisar\_s1}, github.com/mamadou-g-diallo/nisar-s1 (Apache-2.0).
\ifdefined\sansLogo\else Logo UCAD : Wikimedia Commons, CC BY-SA 4.0.\fi\par}
}
\end{columns}
\end{document}
